\begin{center}
    \textbf{ABSTRACT}
\end{center}

\vspace{0.25cm}

%Write a paragraph in English in which a reader could understand what the problem was and what you have done to solve it. Maximum 100 words. You can support this activity by reading the information that you can find in the following hyperlinks: http://www.galaxygoo.org/re-sources/abstract_writing.html http://www.lightbluetouchpaper.org/2007/03/14/how-not-to-write-an-ab-stract/
Inventory replenishment depends on demand forecasts, yet few organizations know which model family best fits their demand. This project developed PRED, an experimental framework that characterizes demand series, trains statistical, machine learning, deep learning and foundation models, compares them through walk-forward validation, Diebold-Mariano tests and a multi-criteria selection protocol, and verifies the winning model with backtesting on unseen demand. It was evaluated on a public Kaggle dataset. [FALTA CERRAR CON ALGUN RESULTADO VISTO]





\begin{center}
    {\textbf{RESUMEN}}
\end{center}

\vspace{0.25cm}

%Redacte un párrafo (máximo 100 palabras) que permita al lector com-prender con claridad cuál era el problema abordado y qué acciones realizó para resolverlo. Esta sección corresponde a la versión en español del abstract y debe mantener la misma precisión y síntesis del texto en inglés.  
La reposición de inventarios depende del pronóstico de demanda, pero pocas organizaciones saben qué familia de modelos se ajusta mejor a su demanda. Este proyecto desarrolló PRED, un framework experimental que caracteriza las series de demanda , entrena modelos estadísticos, de machine learning, de aprendizaje profundo y fundacionales, los compara mediante validación walk-forward, pruebas de Diebold-Mariano y un protocolo de selección multicriterio, y verifica el modelo ganador con backtesting sobre demanda no vista. Se evaluó sobre un conjunto de datos público de Kaggle. [FALTA CERRAR CON ALGUN RESULTADO VISTO]